\documentclass[10pt,conference]{IEEEtran}
\usepackage{amsmath,amssymb}
\usepackage{booktabs}
\usepackage{tabularx}
\usepackage[shortlabels]{enumitem}
\setlist[itemize]{leftmargin=1.4em, itemsep=1pt, topsep=2pt, parsep=0pt}

\newcommand{\st}{\text{s.t.}}
\newcommand{\minimize}{\operatorname*{minimize}}
\newcommand{\maximize}{\operatorname*{maximize}}

\begin{document}

\title{Electricity Market Clearing Problems:\\ A Reference Formulation}

\author{
\IEEEauthorblockN{SOC-Based Storage Market Thesis --- Reference Note}
\IEEEauthorblockA{Technical University of Denmark (DTU)}
}

\maketitle

\begin{abstract}
This note collects, in one consistent notation, the optimization problems that a modern
ISO/RTO solves at each layer of a wholesale electricity market: the day-ahead
security-constrained unit commitment, the real-time security-constrained economic dispatch,
locational marginal pricing, ancillary-service co-optimization, and capacity/resource-adequacy
accreditation. It closes by placing the state-of-charge (SOC) based storage participation
model of Baldick (2026) within this taxonomy, using the same notation as the accompanying
code framework (\texttt{core/}, \texttt{markets/day\_ahead.py}).
\end{abstract}

\begin{IEEEkeywords}
electricity markets, unit commitment, economic dispatch, locational marginal pricing,
ancillary services, resource adequacy, energy storage
\end{IEEEkeywords}

\section{Two-Settlement Structure}
Almost every US and European wholesale market is a \emph{two-settlement system}. A
\textbf{day-ahead (DA) market} clears a financially-binding schedule from forecasts and
offers; a \textbf{real-time (RT) market} then physically balances the system every few
minutes and settles any deviation from the DA schedule at the real-time price. A
\textbf{capacity / resource-adequacy market} sits above both, on a years-ahead horizon --
it prices the obligation to be available, but does not dispatch anything itself. Both DA
and RT jointly clear energy and ancillary services.

\section{Day-Ahead Energy Market}
The day-ahead market solves a security-constrained unit commitment (SCUC): which units
are on, $u_{i,t}\in\{0,1\}$, and their output $g_{i,t}$, for every interval $t\in\mathcal T$.

\begin{align}
\minimize_{u,g,d} \quad &
\sum_{t}\sum_{i}\!\big[c_i(g_{i,t}) + SU_i u_{i,t} + SD_i v_{i,t}\big] \notag\\
&\hspace{2em} - \sum_{t}\sum_{m} u_m(d_{m,t}) \label{eq:scuc-obj}
\end{align}
\begin{align*}
\st \quad
&\sum_i g_{i,t} + \sum_e (g_{e,t}-p_{e,t}) = \sum_m d_{m,t}, && \forall t \tag{a}\\
&u_{i,t} P_i^{\min} \le g_{i,t} \le u_{i,t} P_i^{\max}, && \forall i,t \tag{b}\\
&|g_{i,t}-g_{i,t-1}| \le R_i, && \forall i,t \tag{c}\\
&\text{min up/down-time logic on } u_{i,t}, && \forall i \tag{d}\\
&|f_{\ell,t}| \le \bar f_\ell,\quad f_{\ell,t}=\textstyle\sum_n \Phi_{\ell,n}\,\mathrm{inj}_{n,t}, && \forall \ell,t \tag{e}\\
&\textstyle\sum_i r_{i,t} \ge R_t^{\text{req}}, && \forall t \tag{f}\\
&\text{storage SOC dynamics/limits, if included} \tag{g}
\end{align*}
Constraint (a) is the nodal power balance; (b)--(c) gate output by commitment and ramp
capability; (d) is the logical minimum-run/down-time restriction; (e) is the DC optimal
power flow (DC-OPF) network block; (f)--(g) couple the energy problem to reserves and
storage.

\section{Real-Time Energy Market}
The real-time market re-dispatches continuous output every $\sim\!5$ minutes against the
\emph{actual} net load via security-constrained economic dispatch (SCED); commitment
$u_{i,t}$ is already fixed from the DA solution (or a short-horizon RT commitment for
fast-start units).
\begin{equation}
\minimize_{g} \quad \sum_i c_i(g_{i,t})
\end{equation}
subject to the same balance/capacity/network constraints as \eqref{eq:scuc-obj}, evaluated
at the real-time net load, with two differences: ramp limits are measured from the
\emph{last actual dispatch point} rather than the DA schedule, and any reserve awarded in
the DA market must remain deliverable within its contracted response time.

\section{Locational Marginal Pricing}
The locational marginal price (LMP) is the dual of the nodal power-balance constraint
(a); it settles both the DA and RT energy markets. Whenever the network is modeled
explicitly it decomposes cleanly:
\begin{equation}
\mathrm{LMP}_{n,t} = \lambda_t - \sum_{\ell} \Phi_{\ell,n}\big(\bar\mu_{\ell,t}-\mu_{\ell,t}\big)
\label{eq:lmp}
\end{equation}
where $\lambda_t$ is the reference-bus energy price (common to every node) and
$\bar\mu_{\ell,t},\mu_{\ell,t}\ge 0$ are the congestion duals on line $\ell$'s flow limit,
weighted by the shift factor $\Phi_{\ell,n}$ from $\ell$ to bus $n$. A loss component adds
a third, often-linearized term when losses are modeled explicitly.

\section{Ancillary Services and Reserve Co-Optimization}
\label{sec:as}
Ancillary services are co-optimized jointly with energy in both the DA and RT markets: a
unit that reserves capacity for a product cannot simultaneously sell that same MW as
energy.
\begin{align}
\minimize_{g,r} \quad & \sum_i c_i(g_{i,t}) + \sum_i\sum_k c_i^k(r_{i,t}^k)
\end{align}
with $k \in \{\text{regulation, spinning, non-spinning, ramping}\}$, subject to
\begin{align*}
&\textstyle\sum_i r_{i,t}^k \ge R_t^k, && \forall k,t \tag{a}\\
&g_{i,t}+r_{i,t}^{\text{up}} \le P_i^{\max},\quad g_{i,t}-r_{i,t}^{\text{down}} \ge P_i^{\min}, && \forall i,t \tag{b}\\
&r_{i,t}^k \le \mathrm{RampRate}_i \times \tau_k, && \forall i,k,t \tag{c}\\
&s_{e,t-1} \ge \mathrm{dur}_k\, r_{e,t}^{\text{up}}, \label{eq:soc-reserve}\\
&\bar s_e - s_{e,t-1} \ge \mathrm{dur}_k\, r_{e,t}^{\text{down}}, && \forall e,k,t \tag{d}
\end{align*}
Constraint~(d) is storage-specific: enough energy (respectively headroom) must already
be banked in the SOC to sustain the reserve call for its required duration
$\mathrm{dur}_k$. \textbf{Regulation} (seconds--minutes, AGC-driven, near-zero net energy
over a cycle) is storage's best operational fit; \textbf{spinning/non-spinning reserve}
($\sim\!10$-minute response, sustained 30 minutes to a few hours) is comfortable for most
storage durations; \textbf{voltage support} needs almost no real energy at all.

\section{Capacity / Resource Adequacy Market}
A separate, longer-horizon mechanism: it pays resources to be \emph{available} during
the system's riskiest hours of a future delivery year and does not dispatch anything.
\begin{equation}
\maximize_{\mathrm{Cap}} \quad \sum_i \Big[\mathrm{VRR}(\mathrm{Cap}_i) - \mathrm{Price}_i \cdot \mathrm{Cap}_i\Big]
\end{equation}
cleared against an administrative, downward-sloping demand curve (a Variable Resource
Requirement, or VRR, curve in PJM-style designs), subject to
\begin{align*}
&\textstyle\sum_i \mathrm{ELCC}_i \cdot \mathrm{Cap}_i \ge (1+m)\cdot\mathrm{PeakLoad}, \tag{a}
\end{align*}
where $m$ is the target reserve margin, often set from a Loss-of-Load-Expectation (LOLE)
standard of $\le 0.1$ day/year. $\mathrm{ELCC}_i$ (Effective Load Carrying Capability) is
duration- and penetration-dependent for storage: a 4-hour battery is derated relative to
an 8-hour one, and its \emph{marginal} ELCC falls as more storage of the same duration is
added to the system. Crucially, this is \emph{not} a pass/fail sustained-output test but a
probabilistic, Monte-Carlo loss-of-load measure -- not a daily discharge requirement.

\section{Storage's SOC-Based Participation Model}
\label{sec:storage}
The design this framework reproduces (Baldick, 2026) removes the storage owner's
in-horizon energy offer entirely: the ISO manages the SOC trajectory directly from the
storage's own physical characteristics, priced against a single end-of-horizon bid.
\begin{equation}
\sum_{e,t}\big[\varepsilon_e p_{e,t} + \zeta_e g_{e,t}\big] - \sum_e b_e\big(s_{e,T}-\bar s_{e,T}\big)
\end{equation}
is added to the day-ahead objective~\eqref{eq:scuc-obj}, subject to
\begin{align*}
&s_{e,t} = \gamma_e s_{e,t-1} + \Delta t\,\alpha_e p_{e,t} - \Delta t\,\beta_e g_{e,t}, && \forall e,t \tag{a}\\
&\gamma_e s_{e,t-1}+\Delta t\,\alpha_e p_{e,t} \le \bar s_e, && \forall e,t \tag{b}\\
&\gamma_e s_{e,t-1}-\Delta t\,\beta_e g_{e,t} \ge \underline s_e, && \forall e,t \tag{c}\\
&0\le p_{e,t}\le \bar p_e,\quad 0\le g_{e,t}\le \bar g_e, && \forall e,t \tag{d}
\end{align*}
Constraints (b)--(c) are the \emph{tightened} SOC bounds: they bind the intra-interval
trajectory, not just its endpoints, which is strictly tighter than bounding $s_{e,t}$
alone. The storage owner's opportunity cost is never submitted as a bid -- it is
recovered by the ISO as the dual of the SOC dynamics (a), $\pi_{e,t}$, propagated
backward through the Bellman-like recursion
\begin{equation}
\pi_{e,t} = \gamma_e\,\pi_{e,t+1} - \bar\nu_{e,t+1} + \nu_{e,t+1},
\end{equation}
terminated by $\pi_{e,T} = w_e(s_{e,T}-\bar s_{e,T})$ for a differentiable bid.

\section{Glossary}
\begin{table}[h]
\centering
\small
\begin{tabularx}{\linewidth}{@{}lX@{}}
\toprule
\textbf{Term} & \textbf{Meaning} \\
\midrule
LMP & Locational marginal price (\S IV) \\
SCUC / SCED & Sec.-constrained unit commitment / economic dispatch \\
DC-OPF & DC (linearized) optimal power flow \\
AGC & Automatic Generation Control (regulation signal) \\
ELCC & Effective Load Carrying Capability (\S VI) \\
VRR & Variable Resource Requirement (sloped capacity curve) \\
LOLE & Loss-of-Load Expectation (reliability standard) \\
ISO/RTO & Independent System Operator / Regional Transmission Org. \\
SOC & State of charge (\S VII) \\
Order 841 & FERC's 2018 mandate opening markets to storage \\
\bottomrule
\end{tabularx}
\end{table}

\vfill
\noindent\footnotesize\textit{Companion to \texttt{Code/ROADMAP.md} and
\texttt{Literature/KKT\_MCP\_equilibrium\_SOC\_storage.pdf} -- notation matches
\texttt{core/}, \texttt{markets/day\_ahead.py}.}

\end{document}
